\scanpage{38}
\mspage{30}
% L1–5
Justification a priori de cette description, on considère le graphe ainsi construit à l'aide de $\Delta,V^*$ ; on vérifie qu'il satisfait aux propriétés caractéristiques d'un complexe cubique (th. p. 1)\ldots

\medskip\centerline{\rule{1.2cm}{.3pt}}\medskip
% L6–13
Bases orthogonales à un bitriangle $\Delta$.
Ce sont aussi les bases $b\subset C-\Delta$ (cf p. \unsure{29|24} bis), donc les systèmes de 6 \unsure{élémt|éléments} mutuellmt non liés de $\unsure{F'|T'|\mathbb F'}\times V^*$. Considérons le \unsure{tor.|rev.|tors.} \inserted{$\widetilde I$} d'ordre 2 de $I$ des permutations circulaires sur les fibres de $T\to I$, et considérons $J=\underset{\substack{(\text{produit}\\\text{de Baer}\\i\in I)}}{\bigwedge^{\unsure{\prime}}}\widetilde I_i$, qui est un ens à 2 él. L'ens des bases cherchées est en corr. 1--1 avec $\operatorname{Mon}(J,V^*)$.
% L14–23
Si on a une injection $J\overset{\varphi}{\hookrightarrow}V^*$, on lui associe l'ens \inserted{$b_\varphi$} des $\{x=(\Delta(x),v(x))\in F'\times V^*$ tels que $\varphi(\deleted{\unsure{\Delta}}\text{image de }\Delta(x)\text{ dans }J)=v(x)$ ; ils sont évidemt de cardinal 6 $[\inserted{\text{car }}b_\varphi\xrightarrow{\sim}F')$ \deleted{\unsure{est}}, et ses éléments sont mutuellmt non liés car $v(x)=v(y)$ implique que $\Delta(x)$ et $\Delta(y)$, en \unsure{tant} que $T_i\xrightarrow{\sim}T_j$ $(I=\{i,j\})$, diffèrent par un él. du gpe \uline{alterné} $A_3$, donc $\Delta(x)\cap\Delta(y)=\varnothing$, ce qui implique $x,y$ non liés.
% archive sheet number: 37
